\documentclass[10pt,a4paper]{amsart}
\usepackage{amsmath,amssymb,amsthm,graphicx}
\usepackage{verbatim}
\newenvironment{reference}{\comment}{\endcomment}
%\newenvironment{reference}{}{}
\newenvironment{slogan}{\comment}{\endcomment}
\newenvironment{history}{\comment}{\endcomment}

% For commutative diagrams we use Xy-pic
\usepackage[all]{xy}

% We use 2cell for 2-commutative diagrams.
\xyoption{2cell}
% Theorem environments.
%
\theoremstyle{plain}
\newtheorem{theorem}[subsection]{Theorem}
\newtheorem{proposition}[subsection]{Proposition}
\newtheorem{lemma}[subsection]{Lemma}

\theoremstyle{definition}
\newtheorem{definition}[subsection]{Definition}
\newtheorem{example}[subsection]{Example}
\newtheorem{exercise}[subsection]{Exercise}
\newtheorem{situation}[subsection]{Situation}

\theoremstyle{remark}
\newtheorem{remark}[subsection]{Remark}
\newtheorem{remarks}[subsection]{Remarks}

\numberwithin{equation}{subsection}

% Macros
%
\def\lim{\mathop{\mathrm{lim}}\nolimits}
\def\colim{\mathop{\mathrm{colim}}\nolimits}
\def\Spec{\mathop{\mathrm{Spec}}}
\def\Hom{\mathop{\mathrm{Hom}}\nolimits}
\def\Ext{\mathop{\mathrm{Ext}}\nolimits}
\def\SheafHom{\mathop{\mathcal{H}\!\mathit{om}}\nolimits}
\def\SheafExt{\mathop{\mathcal{E}\!\mathit{xt}}\nolimits}
\def\Sch{\mathit{Sch}}
\def\Mor{\mathop{\mathrm{Mor}}\nolimits}
\def\Ob{\mathop{\mathrm{Ob}}\nolimits}
\def\Sh{\mathop{\mathit{Sh}}\nolimits}
\def\NL{\mathop{N\!L}\nolimits}
\def\CH{\mathop{\mathrm{CH}}\nolimits}
\def\proetale{{pro\text{-}\acute{e}tale}}
\def\etale{{\acute{e}tale}}
\def\QCoh{\mathit{QCoh}}
\def\Ker{\mathop{\mathrm{Ker}}}
\def\Im{\mathop{\mathrm{Im}}}
\def\Coker{\mathop{\mathrm{Coker}}}
\def\Coim{\mathop{\mathrm{Coim}}}

% Boxtimes
%
\DeclareMathSymbol{\boxtimes}{\mathbin}{AMSa}{"02}

%
% Macros for moduli stacks/spaces
%
\def\QCohstack{\mathcal{QC}\!\mathit{oh}}
\def\Cohstack{\mathcal{C}\!\mathit{oh}}
\def\Spacesstack{\mathcal{S}\!\mathit{paces}}
\def\Quotfunctor{\mathrm{Quot}}
\def\Hilbfunctor{\mathrm{Hilb}}
\def\Curvesstack{\mathcal{C}\!\mathit{urves}}
\def\Polarizedstack{\mathcal{P}\!\mathit{olarized}}
\def\Complexesstack{\mathcal{C}\!\mathit{omplexes}}
% \Pic is the operator that assigns to X its picard group, usage \Pic(X)
% \Picardstack_{X/B} denotes the Picard stack of X over B
% \Picardfunctor_{X/B} denotes the Picard functor of X over B
\def\Pic{\mathop{\mathrm{Pic}}\nolimits}
\def\Picardstack{\mathcal{P}\!\mathit{ic}}
\def\Picardfunctor{\mathrm{Pic}}
\def\Deformationcategory{\mathcal{D}\!\mathit{ef}}
\usepackage{cleveref}
\crefname{originalreference}{original reference}{original references}
\Crefname{originalreference}{Original reference}{Original references}
\makeatletter
\newlabel{situation-setup}{{02QL}{source}{}{sourcepointer-1}{}}
\newlabel{situation-setup@cref}{{[originalreference][0][]02QL}{[1][0][]source}{}{}{}}
\makeatother
\crefname{originallemma}{original lemma}{original lemmas}
\Crefname{originallemma}{Original lemma}{Original lemmas}
\makeatletter
\newlabel{more-morphisms-lemma-composition-lci}{{069J}{source}{}{sourcepointer-2}{}}
\newlabel{more-morphisms-lemma-composition-lci@cref}{{[originallemma][0][]069J}{[1][0][]source}{}{}{}}
\makeatother
\crefname{originallemma}{original lemma}{original lemmas}
\Crefname{originallemma}{Original lemma}{Original lemmas}
\makeatletter
\newlabel{lemma-composition-regular-immersion}{{0FF0}{source}{}{sourcepointer-3}{}}
\newlabel{lemma-composition-regular-immersion@cref}{{[originallemma][0][]0FF0}{[1][0][]source}{}{}{}}
\makeatother
\crefname{originallemma}{original lemma}{original lemmas}
\Crefname{originallemma}{Original lemma}{Original lemmas}
\makeatletter
\newlabel{lemma-compose-flat-pullback}{{02RD}{source}{}{sourcepointer-4}{}}
\newlabel{lemma-compose-flat-pullback@cref}{{[originallemma][0][]02RD}{[1][0][]source}{}{}{}}
\makeatother
\title{The Gysin maps of composable local complete intersection morphisms compose}
\author{The Stacks Project Authors}
\date{stacks-project:chow.tex:local-snapshot-20261012; GNU FDL 1.2 or later}
\begin{document}
\maketitle
\noindent\textit{Editorial note.} All delivered mathematical text is the author's own, copied from the pinned Stacks Project chapter file chow.tex; no statement or proof step was rewritten, reordered or removed.\par\medskip
\noindent\textit{Editorial note.} Preamble provenance: the pinned chapter file opens with an input of the shared preamble and carries no preamble text of its own, so the project's single Stacks-wide preamble staged at sources/stacks-chow/preamble.tex is excerpted verbatim as two spans, lines 14-18 (Xy-pic setup) and lines 154-217 (the native theorem declarations and mathematical macros). Nothing was added, removed or redefined; the wrapper class and the amsmath/amssymb/amsthm/graphicx base packages are supplied by the assembler.\par\medskip
\noindent\textit{Editorial note.} The counted statement and proof contain no cross reference, no citation and no included file, so the delivered body is self contained.\par\medskip
\noindent\textit{Editorial note.} Canonical location of the counted proposition: Stacks tag 0FF5 (sources/stacks-chow/pinned-original-source.tex line 13575).\par\medskip
\noindent\textit{Allowed standard prerequisites (editorial).}\begin{itemize}
\item Composition of local complete intersection morphisms. More on Morphisms, Lemma labelled more-morphisms-lemma-composition-lci, cited at proof lines 13583-13585. Makes the statement meaningful.
\item The composition formula for the Gysin maps of regular immersions. Same chapter, Lemma labelled lemma-composition-regular-immersion, cited at proof lines 13608-13609. Supplies the composition of the immersion parts.
\item Composition of flat pullbacks. Same chapter, Lemma labelled lemma-compose-flat-pullback, cited at proof lines 13608-13609. Supplies the composition of the smooth parts.
\item Existence of the Gysin map through an immersion into a smooth scheme. Same chapter, Definition labelled definition-lci-gysin, used at proof lines 13587-13590. Proves that f\textasciicircum{}! exists.
\end{itemize}\medskip
\ifdefined\hypertarget\hypertarget{sourcepointer-1}{}\fi
\noindent\textit{Source reference.} Original reference 02QL: sources/stacks-chow/pinned-original-source.tex, source lines 1372--1372. Same-chapter pointer into the author environment situation at source line 1372; invoked at proof line None of the counted proof. It is retained only as a navigation pointer; every load-bearing step is inside the delivered span.\par\medskip
\ifdefined\hypertarget\hypertarget{sourcepointer-2}{}\fi
\noindent\textit{Source reference.} Original lemma 069J: sources/stacks-chow/tags.txt, source lines 8043--8043. Cross-chapter standard prerequisite located in the project tag list staged in the local snapshot; invoked at proof line 13585 of the counted proof. It is retained only as a navigation pointer; every load-bearing step is inside the delivered span.\par\medskip
\ifdefined\hypertarget\hypertarget{sourcepointer-3}{}\fi
\noindent\textit{Source reference.} Original lemma 0FF0: sources/stacks-chow/pinned-original-source.tex, source lines 13412--13412. Same-chapter pointer into the author environment lemma at source line 13412; invoked at proof line 13606 of the counted proof. It is retained only as a navigation pointer; every load-bearing step is inside the delivered span.\par\medskip
\ifdefined\hypertarget\hypertarget{sourcepointer-4}{}\fi
\noindent\textit{Source reference.} Original lemma 02RD: sources/stacks-chow/pinned-original-source.tex, source lines 2287--2287. Same-chapter pointer into the author environment lemma at source line 2287; invoked at proof line 13607 of the counted proof. It is retained only as a navigation pointer; every load-bearing step is inside the delivered span.\par\medskip
\begin{lemma}
\label{lemma-lci-gysin-composition}
Let $(S, \delta)$ be as in Situation \ref{situation-setup}.
Let $f : X \to Y$ and $g : Y \to Z$ be local complete intersection morphisms
of schemes locally of finite type over $S$. Assume the Gysin
map exists for $g \circ f$ and $g$. Then the Gysin map exists for $f$
and $(g \circ f)^! = f^! \circ g^!$.
\end{lemma}

\begin{proof}
Observe that $g \circ f$ is a local complete intersection morphism
by More on Morphisms, Lemma \ref{more-morphisms-lemma-composition-lci}
and hence the statement of the lemma makes sense.
If $X \to P$ is an immersion of $X$ into a scheme $P$ smooth over $Z$
then $X \to P \times_Z Y$ is an immersion of $X$ into a scheme smooth
over $Y$. This prove the first assertion of the lemma.
Let $Y \to P'$ be an immersion of $Y$ into a scheme $P'$ smooth over $Z$.
Consider the commutative diagram
$$
\xymatrix{
X \ar[r] \ar[d] &
P \times_Z Y \ar[r]_a \ar[ld]^p &
P \times_Z P' \ar[ld]^q \\
Y \ar[r]_b \ar[d] &
P' \ar[ld] \\
Z
}
$$
Here the horizontal arrows are regular immersions, the south-west arrows
are smooth, and the square is cartesian. Whence
$a^! \circ q^* = p^* \circ b^!$ as bivariant classes commute
with flat pullback. Combining this fact with
Lemmas \ref{lemma-composition-regular-immersion} and
\ref{lemma-compose-flat-pullback}
the reader finds the statement of the lemma holds true.
Small detail omitted.
\end{proof}

\end{document}
